% ── Section: Introduction ─────────────────────────────────────────────────────
\section{Introduction}
\label{sec:intro}

The deployment of transformer-based language models in clinical medicine
has produced a new class of infrastructure challenge: \emph{access-controlled
approximate nearest-neighbor (ANN) search} over protected health information
(PHI).
Consider a hospital AI system that helps clinicians find similar past cases.
The system stores BioBERT~\cite{lee2020biobert} embeddings of one million
clinical notes (768-dimensional float vectors) and must answer each query
in under 10 milliseconds.
Crucially, every query result must satisfy the Health Insurance Portability
and Accountability Act (HIPAA): an Oncology physician must not see Psychiatry
notes; a researcher may access only records with explicit consent for their
specific study.

\smallskip
\noindent\textbf{The selectivity problem.}
In Role-Based Access Control (RBAC)~\cite{ferraiolo2001proposed}, each user
holds a \emph{permission mask} that determines which records they may read.
The \emph{selectivity} of a query is the fraction of the database accessible
to the querying user.
In a realistic hospital, selectivity varies from $\approx\!80\%$ (senior
administrator with broad access) down to $\approx\!0.1\%$ (researcher querying
genomic data from HIV-positive trial participants).
No existing ANN index handles this full range efficiently:

\begin{itemize}[leftmargin=1.4em,itemsep=1pt,topsep=2pt]
  \item \textbf{Post-filtering.}
    Standard HNSW retrieves a candidate set and filters by access mask.
    At 0.1\% selectivity, only $\approx\!1{,}000$ of $1{,}000{,}000$ records
    are accessible; a beam of width $\mathit{ef}\!=\!200$ explores far too
    few nodes to reach 10 accessible results reliably---recall collapses
    toward zero~\cite{malkov2020efficient}.

  \item \textbf{Pre-filtering (brute-force).}
    Build the accessible subset on-the-fly and scan it exhaustively.
    Exact results, but $O(|\text{accessible}| \times d)$ work per query:
    at 40\% selectivity over one million vectors, this reduces throughput
    to single-digit queries per second (QPS)~\cite{gollapudi2023filtered}.

  \item \textbf{Multi-index architectures (e.g., SIEVE~\cite{zhang2025sieve}).}
    Maintain one HNSW index per role.
    Good recall and throughput, but memory scales with $|\text{roles}|$:
    64 roles $\Rightarrow$ 64$\times$ the baseline index footprint
    ($\approx\!213$~GB for one million 768-dim vectors; Section~\ref{sec:eval}).
\end{itemize}

\smallskip
\noindent\textbf{Our contribution: RBAC-HNSW.}
We present a \emph{single algorithmic modification} to the HNSW greedy
search that gates on a 64-bit bitmask before computing the vector distance.
Denied nodes are not pruned from graph traversal---their neighbor lists are
still used to \emph{route} the beam toward accessible regions---preventing
the ``connectivity deserts'' that cause post-filtering recall to collapse.

Concretely, the access check
$(\mathit{node\_mask} \;\&\; \mathit{query\_mask}) = \mathit{query\_mask}$
costs one CPU instruction ($\approx\!1$~ns) and is evaluated before the
$768$-multiply cosine distance ($\approx\!30$~ns).
The mask is stored inline with each HNSW node pointer, fitting both in a
single 64-byte CPU cache line~\cite{drepper2007every}.
Memory overhead: $8$~bytes per vector $=$ \textbf{0.24\%} above vanilla HNSW.

\smallskip
\noindent\textbf{Scope of this paper.}
This work targets the \emph{data management community}, framing access-controlled
vector search as a database systems problem.
We present the algorithm (Section~\ref{sec:algorithm}), a biologically motivated
evaluation dataset (Section~\ref{sec:dataset}), and empirical results across
five selectivity levels (Section~\ref{sec:eval}).
Our open-source implementation is available at
\url{https://github.com/patriotsbreeze/Similarity-RBAC}.
